The physical index set is $I=\{0,\ldots,d-1\}$, the bond algebra is
$\MN{D}$, and $G$ denotes a group.  All sums below are finite.  The first
section treats permutation actions on $I$; the second isolates scalar
uniqueness for two gauges of the same injective tensor; the third records the
elementary factor-system bookkeeping used in the cohomology discussion of
Chapter~\ref{ch:symmetry}.

\section{Permutation reindexing and identity bookkeeping}
\label{sec:app_symmetry_permutations}

The linear-combination twist $\widetilde{A}_g^i = \sum_j U(g)_{ij}\, A^j$
is the standard form used in the virtual representation theorem.
The following alternative formulation replaces the matrix action by a
permutation of the physical index; it is useful for symmetries that act by
reordering basis states rather than by a general linear map. % The statements of this section are auxiliary constructions not
% directly sourced from~\cite{PerezGarcia2008StringOrder}.

\begin{definition}[Permutation symmetry datum]\label{def:on_site_symmetry_perm}
    \lean{MPSTensor.OnSiteSymmetry}
    \leanok
    An \emph{on-site symmetry datum} for a group $G$ on a physical
    space of dimension $d$ consists of a matrix-valued map
    $u : G \to \MN{d}$ and a physical-index action
    $\sigma : G \to \mathrm{End}(\{0,\ldots,d{-}1\})$.
\end{definition}

\begin{definition}[Permutation-twisted tensor]\label{def:perm_twisted_tensor}
    \lean{MPSTensor.TwistedTensor}
    \leanok
    \uses{def:on_site_symmetry_perm, def:mps_tensor}
    Given an on-site symmetry datum $(u, \sigma)$, the
    \emph{permutation-twisted tensor} at group element $g$ is
    \begin{align}
        (A^{\sigma_g})^i & := A^{\sigma(g)(i)}.
        \label{eq:symmetry_perm_twist}
    \end{align}
    In tensor-network notation the local tensor is unchanged and only the
    physical leg is relabelled:
    % Formula: (A^{\sigma_g})^i = A^{\sigma(g)(i)}.
    % Source verdict: the adjacent defining identity is the source; there is
    % no dedicated source figure.
    % Ink-to-index: i and \sigma(g)(i) label the free upper physical stub of
    % the same one-site tensor A; the relabelling arrow carries no index.
    % Contractions: none.
    % Paired boundary: boundary|virtual-west=1|virtual-east=1|
    % physical-up=1|physical-down=0.
    \begin{align}
        \begin{tenkz}[physical=up]
            \tn[up=$i$]{A}
        \end{tenkz}
         & \qquad \xrightarrow{\sigma_g} \qquad
        \begin{tenkz}[physical=up]
            \tn[up=$\sigma(g)(i)$]{A}
        \end{tenkz}.
        \notag
    \end{align}
\end{definition}

\begin{lemma}[Identity permutation twist]\label{lem:perm_twisted_one}
    \lean{MPSTensor.TwistedTensor_one}
    \leanok
    \uses{def:perm_twisted_tensor}
    If $\sigma(1) = \id$, then $(A^{\sigma_1})^i = A^i$ for all~$i$.
\end{lemma}

\begin{proof}\leanok\uses{def:perm_twisted_tensor}
    The claim follows from $\sigma(1)(i) = i$ for all~$i$.
\end{proof}

\noindent\emph{Composition order.}
The linear-combination twist (Lemma~\ref{lem:twisted_tensor_mul}) composes as
$\widetilde{(\widetilde{A}_h)}_g = \widetilde{A}_{gh}$, applying $h$ first
then~$g$. The permutation twist below uses the same left-action convention
$\sigma(gh) = \sigma(g) \circ \sigma(h)$, so
$(A^{\sigma_g})^{\sigma_h} = A^{\sigma_{gh}}$; the inner permutation
$\sigma(h)$ acts first on the index.

\begin{lemma}[Composition of permutation twists]\label{lem:perm_twisted_mul}
    \lean{MPSTensor.TwistedTensor_mul}
    \leanok
    \uses{def:perm_twisted_tensor}
    If $\sigma(gh) = \sigma(g) \circ \sigma(h)$ for all $g, h \in G$,
    then
    \begin{align}
        A^{\sigma_{gh}} & = (A^{\sigma_g})^{\sigma_h}.
        \label{eq:symmetry_perm_composition}
    \end{align}
\end{lemma}

\begin{proof}\leanok\uses{def:perm_twisted_tensor}
    Expanding both sides gives
    $(A^{\sigma_{gh}})^i
        = A^{\sigma(g)(\sigma(h)(i))}
        = ((A^{\sigma_g})^{\sigma_h})^i$.
    Diagrammatically this is just two consecutive relabellings of the same
    local tensor:
    % Formula: A^i \xrightarrow{\sigma_h} A^{\sigma(h)(i)}
    %   \xrightarrow{\sigma_g} A^{\sigma(g)(\sigma(h)(i))}.
    % Source verdict: the adjacent componentwise equality is the source;
    % there is no dedicated source figure.
    % Ink-to-index: the three upper stubs carry, from left to right, the free
    % physical indices i, \sigma(h)(i), and \sigma(g)(\sigma(h)(i)).  The two
    % arrows record relabelling by \sigma_h and then \sigma_g; they are not
    % tensor-network contractions.
    % Contractions: none in any panel.
    % Common boundary: boundary|virtual-west=1|virtual-east=1|
    % physical-up=1|physical-down=0.
    \begin{align}
        \tnpic[physical=up]{\tn[up=$i$]{A}}
         & \qquad \xrightarrow{\sigma_h} \qquad
        \tnpic[physical=up]{\tn[up=$\sigma(h)(i)$]{A}}
        \qquad \xrightarrow{\sigma_g} \qquad
        \tnpic[physical=up]{\tn[up=$\sigma(g)(\sigma(h)(i))$]{A}}.
        \notag
    \end{align}
\end{proof}



\subsection{Linear physical-index rotations}

The same bookkeeping applies to a general matrix acting on the physical index.

\begin{definition}[Physical-index rotation]%
    \label{def:rotate_physical}
    \lean{MPSTensor.rotatePhysical}
    \leanok
    \uses{def:mps_tensor}
    Given a matrix $M \in \MN{d}$ and an MPS tensor~$A$,
    the \emph{physical-index rotation} is the tensor
    $(A_M)^i := \sum_j M_{ij}\, A^j$.
    When the matrix is a unitary $u$, the rotation is written $A_u$.
\end{definition}

\begin{lemma}[Matrix product vector under physical rotation]
    \label{lem:mpv_rotate_physical}
    \lean{MPSTensor.mpv_rotatePhysical}
    \leanok
    \uses{def:rotate_physical, def:mpv}
    For a length-$N$ physical configuration
    $s=(s_0,\ldots,s_{N-1})$, physical rotation satisfies
    \begin{align}
        \tr((A_M)^{s_0}\cdots(A_M)^{s_{N-1}})
         & =
        \sum_{t\in\{0,\ldots,d-1\}^N}
        \left(\prod_{n=0}^{N-1}M_{s_n,t_n}\right)
        \tr(A^{t_0}\cdots A^{t_{N-1}}).
        \label{eq:symmetry_mpv_rotation}
    \end{align}
\end{lemma}

\begin{proof}\leanok
    Substitute $(A_M)^{s_n}=\sum_{t_n}M_{s_n,t_n}A^{t_n}$ at each
    site, distribute the finite sums through the ordered matrix product, and
    use linearity of the trace to obtain
    (\ref{eq:symmetry_mpv_rotation}).
\end{proof}

\begin{lemma}[Word expansion for a twisted tensor]
    \label{lem:eval_word_twisted_tensor}
    \lean{MPSTensor.evalWord_twistedTensor_ofFn}
    \leanok
    \uses{def:twisted_tensor}
    For a word $b:\{0,\ldots,L-1\}\to\{0,\ldots,d-1\}$,
    \begin{align}
        A_g^{b_0}\cdots A_g^{b_{L-1}}
         & =\sum_{v:\{0,\ldots,L-1\}\to\{0,\ldots,d-1\}}
        \left(\prod_{k=0}^{L-1}U(g)_{b_kv_k}\right)
        A^{v_0}\cdots A^{v_{L-1}}.
        \label{eq:app_symmetry_twisted_word}
    \end{align}
\end{lemma}

\begin{proof}\leanok
    Substitute $A_g^{b_k}=\sum_{v_k}U(g)_{b_kv_k}A^{v_k}$ into the ordered
    product.  Repeated distributivity gives one summand for each function
    $v$, scalar coefficients commute past the matrices and multiply, and the
    remaining ordered product is $A^{v_0}\cdots A^{v_{L-1}}$.
\end{proof}

\section{Gauge ratios and scalar uniqueness}
\label{sec:app_symmetry_gauge_ratios}

The virtual representation theorem
(Theorem~\ref{thm:virtual_rep_injective}) relies on the fact that
composing two gauge transformations yields another gauge up to a
scalar factor.

\begin{lemma}[Gauge ratio commutes with tensor matrices]%
    \label{lem:gauge_ratio_commutes}
    \lean{MPSTensor.gauge_ratio_commutes}
    \leanok
    \uses{def:gauge_equiv}
    Diagrammatically, the hypothesis is that the two local gauges
    \begin{center}
        % B^i = X A^i X^{-1}: the red capsules act on the west and east
        % virtual indices of A^i; the upper leg is the free index i.
        % Both sides have boundary signature (1,1,1,0).
        \begin{tenkz}[physical=up]
            \tn[up=$i$]{B}
        \end{tenkz}
        $ = $
        \begin{tenkz}[physical=up]
            \tnX{X} & \tn[up=$i$]{A} & \tnX{X^{-1}}
        \end{tenkz}

        \qquad\text{and}\qquad

        % B^i = Y A^i Y^{-1}: this second gauge has the same source,
        % target, and free physical index as the X-gauge above.  Both
        % sides again have boundary signature (1,1,1,0).
        \begin{tenkz}[physical=up]
            \tn[up=$i$]{B}
        \end{tenkz}
        $ = $
        \begin{tenkz}[physical=up]
            \tnX{Y} & \tn[up=$i$]{A} & \tnX{Y^{-1}}
        \end{tenkz}
    \end{center}
    produce the same target tensor $B^i$ from the same source tensor
    $A^i$.
    If $X, Y \in \GL_D(\C)$ both conjugate an MPS tensor $A$ to the same
    tensor $B$ (i.e.\ $B^i = X A^i X^{-1} = Y A^i Y^{-1}$ for all $i$),
    then $Y^{-1} X$ commutes with every $A^i$.
\end{lemma}

\begin{proof}\leanok
    Direct matrix calculation: conjugating both gauge relations and
    composing gives $(Y^{-1} X) A^i = A^i (Y^{-1} X)$.
\end{proof}

\begin{lemma}[Gauge ratio is scalar for injective tensors]%
    \label{lem:gauge_ratio_isScalar}
    \lean{MPSTensor.gauge_ratio_isScalar}
    \leanok
    \uses{lem:gauge_ratio_commutes, def:injective, lem:center_scalar}
    Under the same hypotheses, if $A$ is additionally injective, then
    $Y^{-1} X = \lambda \Id_D$ for some scalar $\lambda \in \C$.
\end{lemma}

\begin{proof}\leanok\uses{lem:gauge_ratio_commutes, lem:center_scalar}
    The same two-gauge picture shows that the ratio $Y^{-1}X$ acts
    invisibly on the local tensor.
    By Lemma~\ref{lem:gauge_ratio_commutes}, $Y^{-1}X$ commutes with
    all $A^i$.
    Since $\{A^i\}$ spans $\MN{D}$, the matrix $Y^{-1}X$ commutes
    with all of $\MN{D}$.
    By Lemma~\ref{lem:center_scalar}, it is a scalar matrix.
\end{proof}

\begin{theorem}[Projective multiplication law for gauge matrices]%
    \label{thm:gaugeMatrix_projective_mul}
    \lean{MPSTensor.gaugeMatrix_projective_mul}
    \leanok
    \uses{lem:gauge_ratio_isScalar, def:injective,
        def:on_site_symmetric, lem:twisted_tensor_mul}
    Let $A$ be an injective MPS tensor with on-site symmetry under a
    group representation $U : G \to \GL_d(\C)$.
    For each $g \in G$, let $X(g) \in \GL_D(\C)$ satisfy
    $\widetilde{A}_g^i = X(g)\, A^i\, X(g)^{-1}$ for all~$i$.
    Then for all $g, h \in G$, there exists $c(g,h) \in \C$ such that
    $X(h)\, X(g) = c(g,h)\, X(g \cdot h)$.
    In fact $c(g,h) \ne 0$ since both sides are invertible, so
    $c(g,h) \in \C^{\times}$. Diagrammatically, the two candidate gauges
    are competing local replacements for the same twisted tensor:
    \begin{center}
        % (X(h)X(g)) A^i (X(h)X(g))^{-1}
        %   = X(gh) A^i X(gh)^{-1}: the two candidate gauges act on
        % the same virtual indices of A^i and produce the same twisted
        % tensor.  The upper leg is i; both terms have signature (1,1,1,0).
        \begin{tenkz}[physical=up]
            \tnX{X(h)X(g)} & \tn[up=$i$]{A} &
            \tnX{(X(h)X(g))^{-1}}
        \end{tenkz}
        $ = $
        \begin{tenkz}[physical=up]
            \tnX{X(gh)} & \tn[up=$i$]{A} & \tnX{X(gh)^{-1}}
        \end{tenkz}
    \end{center}
    and both sides conjugate $A^i$ to the same twisted tensor
    $\widetilde{A}_{gh}^i$.
\end{theorem}

\begin{proof}\leanok
    \uses{lem:gauge_ratio_isScalar, lem:twisted_tensor_mul}
    By the composition law $\widetilde{A}_{gh} = \widetilde{(\widetilde{A}_h)}_g$
    (Lemma~\ref{lem:twisted_tensor_mul}), both $X(h)\,X(g)$ and $X(gh)$
    conjugate $A$ to $\widetilde{A}_{gh}$.
    Lemma~\ref{lem:gauge_ratio_isScalar} then gives the scalar factor.
\end{proof}

\section{Cocycle equivalence and coboundaries}
\label{sec:app_symmetry_cocycles}

Let $\omega_1,\omega_2:G\times G\to\C^\times$ be arbitrary scalar
$2$-cochains.  The definitions of coboundary and cohomology are
Definitions~\ref{def:is_coboundary} and~\ref{def:cohomologous_to}.

\begin{theorem}[Cohomology relation is an equivalence relation]%
    \label{thm:cohomologous_equivalence}
    \lean{TNLean.Algebra.ScalarCocycle.CohomologousTo.equivalence}
    \leanok
    \uses{def:cohomologous_to}
    The relation ``cohomologous to'' is reflexive, symmetric, and
    transitive on the set of $\C^\times$-valued scalar $2$-cochains, and
    therefore also on the subset of cocycles.
\end{theorem}

\begin{proof}\leanok\uses{def:cohomologous_to}
    Reflexivity: take $\varphi=1$.
    Symmetry: replace $\varphi$ by $\varphi^{-1}$.
    Transitivity: compose the witnesses $\varphi$ and $\psi$ pointwise
    in (\ref{eq:symmetry_cohomologous}).
\end{proof}


\begin{lemma}[Coboundary iff cohomologous to the trivial cocycle]%
    \label{lem:coboundary_iff_cohomologous_one}
    \lean{TNLean.Algebra.ScalarCocycle.isCoboundary_iff_cohomologousTo_one}
    \leanok
    \uses{def:is_coboundary, def:cohomologous_to}
    A scalar $2$-cochain $\omega$ is a coboundary if and only if
    $\omega \sim \Id$, where $\Id(g,h)=1$ for all $g,h$.
\end{lemma}

\begin{proof}\leanok\uses{def:is_coboundary, def:cohomologous_to}
    Both directions follow by absorbing or introducing the factor
    $\Id(g,h)=1$ via $x\cdot 1=x$ in
    (\ref{eq:symmetry_cohomologous}).
\end{proof}
